\UseRawInputEncoding

\lstdefinestyle{promptstyle}{
  basicstyle=\ttfamily\scriptsize,
  breaklines=true,
  breakatwhitespace=false,
  columns=fullflexible,
  keepspaces=true,
  showstringspaces=false,
  tabsize=2
}

\begin{table*}[t]
\begin{tcblisting}{
  listing only,
  enhanced,
  colback=lightblue!5,
  colframe=lightblue!50,
  rounded corners,
  top=6pt,bottom=6pt,left=8pt,right=8pt,boxsep=4pt,
  title=Prompt for General Domain Judge Generation ({\dataset} Curation),
  listing options={style=promptstyle}
}
You are a fair and impartial judge. Your task is to evaluate 'Response A' and 'Response B' based on a given instruction and a rubric. You will conduct this evaluation in distinct phases as outlined below.

### Phase 1: Compliance Check Instructions
First, identify the single most important, objective 'Gatekeeper Criterion' from the rubric.
- **A rule is objective (and likely a Gatekeeper) if it can be verified without opinion. Key examples are: word/paragraph limits, required output format (e.g., JSON validity), required/forbidden sections, or forbidden content.**
- **Conversely, a rule is subjective if it requires interpretation or qualitative judgment. Subjective rules about quality are NOT Gatekeepers. Examples include criteria like "be creative," "write clearly," "be engaging," or "use a professional tone."**
Think step-by-step to determine this single most important Gatekeeper, then write a 1-2 sentence explanation of your decision.

### Phase 2: Analyze Each Response
Next, for each Gatekeeper Criterion and all other criteria in the rubric, evaluate each response item by item.
For each item, think step-by-step and cite concrete evidence from the response before assigning your judgment.

### Phase 3: Final Judgment Instructions
Based on the results from the previous phases, determine the winner using these simple rules. Provide a final justification explaining your decision first and then give your decision.
Think step-by-step to aggregate the findings and make the decision; keep the reasoning explicit and concise.

---
### REQUIRED OUTPUT FORMAT
You must follow this exact output format below.

--- Compliance Check ---
Gatekeeper Reasoning: <1-2 sentences citing the relevant rubric text>
Identified Gatekeeper Criterion: <e.g., Criterion 1: Must be under 50 words.>

--- Analysis ---
**Response A:**
- Criterion 1 [Hard Rule]: Justification: <...>
- Criterion 2 [Hard Rule]: Justification: <...>
- Criterion 3 [Principle]: Justification: <...>
- ... (and so on for all other criteria)

**Response B:**
- Criterion 1 [Hard Rule]: Justification: <...>
- Criterion 2 [Hard Rule]: Justification: <...>
- Criterion 3 [Principle]: Justification: <...>
- ... (and so on for all other criteria)

--- Final Judgment ---
Aggregation Summary: <1-3 sentences explaining how Gatekeeper and other criteria led to the decision>
Justification: <...>
Winner: <Response A / Response B>


Task to Evaluate:
Instruction:
{instruction}

Rubric:
{rubric}

Response A:
{response_a}

Response B:
{response_b}
\end{tcblisting}
\end{table*}